\documentclass[11pt]{article}
\usepackage[margin=1in]{geometry}
\usepackage[T1]{fontenc}
\usepackage{lmodern,amsmath,amssymb,amsthm,mathtools,hyperref}
\setlength{\emergencystretch}{3em}
\newtheorem{theorem}{Theorem}[section]
\newtheorem{lemma}[theorem]{Lemma}
\newtheorem{corollary}[theorem]{Corollary}
\newtheorem{proposition}[theorem]{Proposition}
\newcommand{\E}{\mathbb E}
\newcommand{\tr}{\operatorname{tr}}
\title{Continuous phase escape, bounded completion,\\and larger discrete exchanges}
\author{Joshua Mark\thanks{ORCID: \href{https://orcid.org/0009-0003-5905-2382}{0009-0003-5905-2382}. AI tools assisted with derivation, independent review, computation and exposition.}}
\date{20 November 2026 --- research draft; priority not established}
\setcounter{MaxMatrixCols}{20}
\begin{document}\maketitle

\begin{abstract}
Every real sign matrix that is not Hadamard is a strict saddle of
the complex-phase squared-Gram energy. An exact construction supplies
integer phase directions, escaping even the previously certified
discrete one- and two-entry traps. Separately, classical real completion
theory yields quaternary completion with at most three columns missing,
with a bounded candidate enumeration algorithm. We also examine
larger discrete exchanges and support-preserving weighing moves.
These results provide escape and completion tools; they do not produce
the required inputs at every admissible order or prove the conjecture.
Bibliographic priority is unestablished.
\end{abstract}
The real Hadamard conjecture asks for a sign matrix $H$ satisfying
$HH^T=nI$ at every positive order divisible by four.
Moving into complex phases changes the feasible set; returning to a
finite phase alphabet is an additional requirement throughout this report.
\section{Every non-Hadamard real sign matrix is a phase saddle}
For an $n$-square matrix $F$ with unit-modulus complex entries put
\[
 \Psi(F)=\|FF^*-nI\|_F^2.
\]
At a real sign matrix $H$, use phase coordinates
$F(\Theta)_{ia}=h_{ia}e^{i\theta_{ia}}$ and let $L_H$ be the real
Hessian at $\Theta=0$. The first derivatives vanish: every rectangle
term in the energy is a real sign times a cosine.
This phase landscape and its flows are classical background;
see Motta--Shipman, \emph{Informing the structure of complex Hadamard
matrix spaces using a flow},
\url{https://doi.org/10.3934/dcdss.2019147}.
The argument below is independently derived; priority is unestablished.

\begin{theorem}[Strict phase saddle at every nonorthogonal sign matrix]
For $n\ge2$, a real sign matrix $H$ has positive-semidefinite
phase Hessian $L_H$ if and only if $HH^T=nI$.
Otherwise an explicit real phase direction has strictly negative
second derivative. In particular every nonglobal discrete trap
from the previous pass admits infinitesimal complex-phase escape.
\end{theorem}
\begin{proof}
Put $G=HH^T$ and $C=GH$. If $H$ is Hadamard, $\Psi=0$ is a global
minimum, so its Hessian is positive semidefinite.

For any two columns $a,b$, multiply $H$ on the right by the unitary
two-square block
\[
 \begin{pmatrix}\cos t&i\sin t\\i\sin t&\cos t\end{pmatrix}.
\]
Every resulting entry still has modulus one, because its real and
imaginary parts are signed $\cos t,\sin t$. Right multiplication
preserves the row Gram matrix, hence $\Psi$ is exactly constant.
Its phase tangent $v$ has entries
$v_{ia}=v_{ib}=h_{ia}h_{ib}$ and all other entries zero.
Therefore $v^TL_Hv=0$.

The phase gradient is
$4\operatorname{Im}(\overline{F_{ia}}(FF^*F)_{ia})$.
Along this path $FF^*F=CH_{\rm block}(t)$, so its derivative in
column $a$ is $4(h_{ia}C_{ib}-h_{ib}C_{ia})$.
If $L_H$ were positive semidefinite, every zero-quadratic direction
would be in its kernel. Thus, for every $i,a,b$,
$h_{ia}C_{ib}=h_{ib}C_{ia}$, or $C=\Lambda H$ for a real diagonal
matrix $\Lambda$.

Left multiplication by the same two-square block on any pair of
rows likewise preserves unit modulus and $\Psi$. Differentiating
its gradient gives $h_{ia}C_{ja}=h_{ja}C_{ia}$ for all $i,j,a$.
Consequently $C=H M$ for a real diagonal $M$. Since every entry of
$H$ is nonzero, $\Lambda H=HM$ forces both diagonals to be a
common scalar: $GH=\lambda H$.
It follows that every nonzero eigenvalue of $G$ equals $\lambda$.
Writing $r=\operatorname{rank}H$, its trace gives $\lambda=n^2/r>0$.
If $r=n$ this is exactly $G=nI$.

If $r<n$, choose nonzero real $u\in\ker H^T$, $v\in\ker H$,
put $X=uv^T$, and take phase direction $\Theta=H\circ X$.
Its first matrix derivative is $iX$, and its second is
$-H\circ X^{\circ2}$. Both $HX^T$ and $XH^T$ vanish.
The first Gram derivative is therefore zero, while the second is
\[
 2XX^T-(H\circ X^{\circ2})H^T
       -H(H\circ X^{\circ2})^T.
\]
Taking twice its inner product with $G-nI$ yields
\[
 \Theta^TL_H\Theta
 =-4n\|X\|_F^2-4(\lambda-n)\|X\|_F^2
 =-4\lambda\|X\|_F^2<0,
\]
contradicting positive semidefiniteness.
\end{proof}

\subsection{An exact algorithm for the escape direction}
The proof can be made constructive without numerical eigenvalues.
For every row- and column-pair tangent $v$, compute $w=L_Hv$.
The exact constant-energy paths ensure $v^TL_Hv=0$.
If some $w\ne0$, let
$B=\max_i\sum_j|(L_H)_{ij}|$ and use
\[
 y=(B+1)v-w.
\]
Then
$y^TL_Hy=-2(B+1)\|w\|^2+w^TL_Hw
\le-(B+2)\|w\|^2<0$.
If all these products vanish, the proof gives $GH=\lambda H$.
Either $H$ is already Hadamard or rational null vectors supply
the negative direction above. Scaling clears all denominators.
Thus real sign inputs admit integer phase directions with exact
negative-curvature certificates.

The Hessian itself has an integer rectangle formula. For
$i<j,a<b$, let $v_{ijab}$ have $+1$ at $(i,a),(j,b)$ and
$-1$ at $(i,b),(j,a)$. Then
\[
 L_H=-4\sum_{i<j,\ a<b}
 h_{ia}h_{ib}h_{ja}h_{jb}\,v_{ijab}v_{ijab}^T.
\]
The checker compares the construction with all normalized order-four
sign matrices and explicit order-eight and order-sixteen traps.

An actual decreasing step can also be certified with rational
arithmetic. Given the nonzero integer direction $y$, set
$s_{ia}=y_{ia}/\max|y|$ and replace each sign by
\[
 F(t)_{ia}=h_{ia}\frac{1+it s_{ia}/2}{1-it s_{ia}/2}.
\]
For rational $t$, all entries are Gaussian rational and have modulus
one. At zero this path has first and second derivatives $iH\circ s$
and $-H\circ s^{\circ2}$, hence the same strictly negative phase
second derivative. Exact backtracking through $t=1,1/2,1/4,\ldots$,
accepting only $\Psi(F(t))<\Psi(H)$, must terminate. This is a
finite exact escape algorithm, not merely an eigenvalue assertion.

\subsection{What escape does and does not provide}
Every sufficiently short nonzero step in the certified direction
reduces $\Psi$. This defeats the discrete trap by enlarging the
allowed alphabet continuously. It does not prove that the trajectory
converges, that other complex critical points are saddles, or that its
limit has four phases. Complex Hadamard matrices already exist at
every order, including orders where real Hadamard matrices cannot exist.
The missing phase-quantization step is therefore essential.
\section{Three missing quaternary columns: a classical-input completion algorithm}
Let $n$ be even and let $V\in\mu_4^{n\times(n-r)}$ satisfy
$V^*V=nI$, with $1\le r\le3$. The following completion statement
uses the \emph{classical real codimension-at-most-seven Hadamard
completion theorem}. That input is explicitly assumed here, not
reproved or claimed new. References include Hall's four-row result
and Verheiden's extension, recorded in
\url{https://thesis.caltech.edu/2899/}, and the primary discussion in
\emph{Counting results for thin Butson matrices},
\url{https://www.combinatorics.org/ojs/index.php/eljc/article/download/v21i3p12/pdf/}.

\begin{lemma}[A one-column extension gate through realification]
For arbitrary $r\ge1$, a quaternary column orthogonal to all columns
of $V$ exists if and only if the real sign frame
$A=\mathcal R((1+i)V)$ has an orthogonal real sign column.
\end{lemma}
\begin{proof}
Write a real sign column in paired coordinates as
$w=(a_1,b_1,\ldots,a_n,b_n)^T$ and let $u_i=a_i+ib_i$.
The paired realified columns give
$A^Tw=0$ if and only if $((1+i)V)^*u=0$, equivalently $V^*u=0$.
Thus $z=u/(1+i)\in\mu_4^n$ is the desired new column.
Conversely $(1+i)z$ has real and imaginary parts in $\{\pm1\}$,
and realifies to such a $w$.
\end{proof}

\begin{theorem}[Automatic completion in complex codimension at most three]
Every such $V$ with $1\le r\le3$ can be completed to a
$\mathrm{BH}(n,4)$ matrix, retaining all its original columns.
\end{theorem}
\begin{proof}
For even $n\ge4$, $A$ has $2n$ rows, $2n-2r$ sign columns and
$A^TA=2nI$. Its real codimension is $2r\le6$, so classical real
completion supplies a full real Hadamard extension and in particular
one new real sign column orthogonal to $A$. The lemma converts that
column to one quaternary column orthogonal to $V$. Append it and
repeat at complex codimension $r-1$. The order is unchanged.
For $n=2$ the assertion follows directly from the order-two matrix:
one prescribed fourth-root column is completed by changing one sign,
and an empty input uses that matrix itself.
\end{proof}
This avoids an unwarranted assumption that an arbitrary supplied real
completion respects the canonical complex structure. Only one real
column is used at each step; its canonical complex-structure partner
is automatically another sign column. The unused real completing
columns may be discarded, and the classical theorem is applied again.

\begin{theorem}[An explicit finite exact algorithm]
For $r\le3$, Gaussian-rational row reduction and at most
$4^{r-1}+\cdots+4+1\le21$ candidate-column tests construct the
completion. The theorem does not require a real Hadamard completion
as algorithmic input.
\end{theorem}
\begin{proof}
Compute a systematic basis of $\ker V^*$: Gaussian elimination
selects $r$ free coordinate positions, and every kernel vector is
determined by its values there. A fourth-root vector can be globally
multiplied by a fourth root so that its first free coordinate equals
one. Its other free coordinates have $4^{r-1}$ possible assignments.
For each assignment compute the unique kernel vector and check all
its entries are in $\mu_4$. The preceding theorem guarantees a
successful assignment. Append that vector and repeat. Its squared
norm is $n$, and its orthogonality to all existing columns is exact,
so the invariants hold at every stage. All coefficients are Gaussian
rational and every test can be done exactly.
\end{proof}
For fixed $r\le3$ the matrix elimination costs $O(n^3)$ scalar
operations; after a kernel basis is available each candidate costs
$O(nr)$ operations. This is an exact-arithmetic operation count,
not an unrestricted bit-complexity or approximate-input claim.

\begin{corollary}[A precise codimension-four gate]
For a quaternary input with $r=4$, full quaternary completion is
equivalent to the existence of even one new fourth-root column.
The latter is decided by at most $4^3=64$ normalized candidate tests.
If a new column exists, the remaining defect three is automatically
completed by the theorem. No automatic existence of that first
column at defect four is asserted.
\end{corollary}
In particular a maximal proper quaternary frame at even order must
miss at least four columns. The result closes a fixed-size residual
gap once a large partial frame is supplied; it does not construct
that input at unknown orders. The completion theory and its small-
alphabet consequence are not advertised as bibliographically new.
\section{Constructive strip and row escapes from the discrete tensor traps}
\label{sec:exchanges-layer-escape}
The preceding discrete traps obstruct improving one- and two-entry moves.
They do not obstruct larger structured updates. This section gives an
exact decomposition of the energy under Hadamard row layers, a sharp
barrier in the corresponding strip-move graph, and a four-step block-row
repair of the explicit tensor family. It does not provide a repair rule
for arbitrary sign matrices or at unknown Hadamard orders.

Write $F_n(H)=\|HH^T-nI\|_F^2$. Let $B$ be a sign Hadamard matrix of
order $m$ and let $A_1,\ldots,A_m$ be arbitrary order-$N$ sign matrices.
Define the $Nm\times Nm$ sign matrix
\[
 H_{(i,p),(j,q)}=(A_p)_{ij}B_{pq},\qquad
 1\le i,j\le N,\quad1\le p,q\le m.
\]
\begin{theorem}[An exact sum of independent row-layer energies]
\label{thm:exchanges-layer-energy}
After grouping rows by $p$, the row Gram of $H$ is block diagonal with
blocks $mA_pA_p^T$. Consequently
\[
 F_{Nm}(H)=m^2\sum_{p=1}^m F_N(A_p).
\]
In particular $H$ is Hadamard if and only if every $A_p$ is Hadamard.
Flipping one entry of one $A_p$ is realized by flipping an $m$-entry
strip in one actual row of $H$, and its energy change is exactly
$m^2$ times the base entry change. An arbitrary replacement of one base
row is similarly one actual row replacement in $H$.
\end{theorem}
\begin{proof}
The dot product of rows $(i,p)$ and $(i',p')$ is
\[
 \sum_j(A_p)_{ij}(A_{p'})_{i'j}\sum_qB_{pq}B_{p'q}.
\]
It is zero when $p\ne p'$, and equals $m(A_pA_p^T)_{ii'}$ otherwise.
Subtract $NmI$ and square the Frobenius norm blockwise.
The update statements follow directly from the definition.
\end{proof}

Let $A$ be the order-eight matrix with negative-entry row bitmasks
\[
 (84,44,30,230,133,73,111,51),
\]
using zero-based column bits. Its energy is $152$; the preceding pass
certified positive cost for every single-entry move, nonnegative cost
for every two-entry move, and a three-entry escape.
\begin{proposition}[Exact escape barrier for strip moves]
\label{prop:exchanges-strip-barrier}
Starting from $A\otimes B$, the minimum number of strip flips needed
to reach strictly lower energy is three. The least energy barrier along
strip-flip paths to a lower-energy matrix is exactly $8m^2$.
An explicit escape changes only one actual row and $3m$ entries.
These minimum claims concern strip moves, not unrestricted entry moves.
\end{proposition}
\begin{proof}
Every initial strip costs at least $8m^2$. Two strips in different layers
each have positive cost; two in the same layer have nonnegative total
cost by the exact base certificate. For an attaining path choose one
layer and flip base cells $(2,4),(2,2),(2,3)$ in that order. All lie in
one base row. The base energy path is
$152\to160\to152\to128$, so the lifted path is
\[
 152m^3\ \longrightarrow\ 152m^3+8m^2\
 \longrightarrow\ 152m^3\
 \longrightarrow\ 152m^3-24m^2.
\]
It attains both lower bounds.
\end{proof}

\begin{theorem}[Four improving block-row updates repair the whole family]
\label{thm:exchanges-four-row-repair}
For every Hadamard factor $B$, the tensor trap $A\otimes B$ admits an
explicit repair to a Hadamard matrix through four strictly improving
block-row replacements. Their energies are
\[
 152m^3\to120m^3\to96m^3\to56m^3\to0.
\]
Equivalently the repair can be performed by $4m$ strictly improving
ordinary-row replacements. Exactly $7m^2$ entries are changed in this
construction. This count is not asserted to be globally minimal.
\end{theorem}
\begin{proof}
Perform the following base row replacements in order:
\[
 \begin{array}{c|c|c|c}
 \text{row}&\text{old negative mask}&\text{new mask}&\text{energy after update}\\\hline
 3&230&2&120\\
 2&30&31&96\\
 6&111&103&56\\
 7&51&49&0
 \end{array}
\]
The final row masks are
$(84,44,31,2,133,73,103,49)$ and have exact Gram $8I$.
The checker verifies each displayed energy and Gram with integers.
The mask xors have respective Hamming weights four, one, one, one,
so seven base entries change, all at distinct positions.

Apply each update simultaneously to all $m$ layers. The matrix remains
$A^{(t)}\otimes B$, giving the stated $m^3$ energy scaling and final
Hadamard matrix. Each changed base entry affects $m^2$ tensor entries,
so the construction changes $7m^2$ entries. Alternatively update a single
layer at a time: the exact sum formula says each ordinary-row replacement
lowers energy by respectively $32m^2,24m^2,40m^2$, or $56m^2$.
There are $m$ replacements of each of the four kinds and no step
increases energy.
\end{proof}

The independent-layer identity gives a useful search rule when a tensor
or Hadamard-layer structure is known: optimize the small base layers
rather than relying on single-entry descent in the expanded matrix.
It also makes the limit explicit: a Hadamard matrix in this class
requires Hadamard base layers, so the construction cannot prove an
unknown base order by hiding it inside a larger tensor dimension.
Classical block-Hadamard operations are the background; priority for
these particular trap-repair certificates has not been established.
No new Hadamard orders or all-order existence claim is made.
\section{A support-moving four-row repair of a weighing obstruction}
\label{sec:weighing-switch-repair}

The two-zero completion test detects an obstruction but does not
produce a good weighing input. Here a local orthogonal operation
changes such an input while preserving its weighing identity.
The operation moves zero positions as well as signs. On the explicit
unfillable $W(12,10)$ it produces a matrix admitting two Hadamard
fillings. This is a repair at a known order, not an all-order theorem
or a minimum-entry-distance claim.

We use the preceding normalization for $W(n,n-2)$, $4\mid n$:
the zero positions are diagonal two-square blocks and every
off-diagonal block is Hadamard. Orthogonal row mixing and switching
belong to the classical orthogonal-design setting; see, for context,
Geramita--Verner, \emph{Orthogonal Designs with Zero Diagonal},
\url{https://doi.org/10.4153/CJM-1976-028-3}.
The criterion and finite repair below are independently checked;
no novelty claim is made.

\begin{theorem}[A four-row support-moving switch]
Select two complete zero-row pairs, ordered as four rows, and let
\[
 Q=\frac12\begin{pmatrix}
1&1&1&1\\1&-1&1&-1\\1&1&-1&-1\\1&-1&-1&1
\end{pmatrix}.
\]
Replace those four rows by their product with $Q$, leaving all
other rows unchanged. The output is another $W(n,n-2)$ if and
only if the product of the four selected entries is $-1$ in every
column outside the two selected zero-column pairs.
When it succeeds, zero positions can move only inside the four-square
submatrix formed by the selected row and zero-column pairs.
\end{theorem}
\begin{proof}
The row operation is orthogonal, so the weighing Gram identity is
preserved without further conditions. In a selected zero-column pair,
the four-entry column has two zeros and two signs. Multiplication
by $Q$ gives entries in $\{0,\pm1\}$ and squared norm two, hence
exactly two nonzero signs. In any other column all four selected
entries are signs. If their product is $-1$, each Walsh row has
inner product $+2$ or $-2$ with that sign vector, so the output
consists of four signs. If their product is $+1$, each inner product
is zero or $\pm4$; norm preservation forces one output entry of
magnitude two, which is forbidden. This proves necessity and
sufficiency, including the support statement. Since the square output
has Gram $(n-2)I$, its rows and columns still contain exactly two zeros.
\end{proof}

Unlike row permutations and sign changes, this operation need not
preserve zero-fill feasibility. Its support motion is essential to
the following small-support rigidity statement.

\begin{proposition}[Two fixed-support row blocks cannot change feasibility]
Suppose $W'$ is another normalized $W(n,n-2)$ with exactly the same
zero positions as $W$, and only one or two zero-row pairs are allowed
to differ. Then $W'$ is obtained from $W$ by signed row permutations
inside those pairs. In particular their zero-preserving Hadamard
completion problems are equivalent.
\end{proposition}
\begin{proof}
The matrix $U=W'W^{\mathsf T}/(n-2)$ is orthogonal and satisfies
$W'=UW$. Every unchanged row forces the corresponding row and
column of $U$ to be a coordinate vector, so $U$ acts only on the
selected pairs. For two pairs its four-square part has two-square
blocks $\left(\begin{smallmatrix}P&Q_1\\R&P_1\end{smallmatrix}\right)$.
Preserving the zero block in the first pair gives
$Q_1W_{ji}=0$; preserving the second gives $RW_{ij}=0$.
Both off-diagonal blocks are invertible Hadamard matrices, so
$Q_1=R=0$. The remaining blocks $P,P_1$ are orthogonal. Each maps
an off-diagonal two-square Hadamard block to another such block,
and hence is a signed permutation: its entries are row inner products
of two-square Hadamard matrices divided by two, thus $0,\pm1$.
The one-pair case is the same argument without the cross blocks.
Signed row permutations give bijections between sign completions.
\end{proof}

\subsection{An explicit escape from the order-twelve obstruction}
For a self-contained input, put
\[
 F=\begin{pmatrix}1&1\\1&-1\end{pmatrix},\quad
 I=I_2,\quad D=\begin{pmatrix}1&0\\0&-1\end{pmatrix},\quad
 S=\begin{pmatrix}0&1\\1&0\end{pmatrix},
\]
and define the six-square block matrix
\[
 A=\begin{pmatrix}
0&I&D&D&D&D\\S&0&-I&-I&I&I\\S&-D&0&I&-D&D\\
S&-D&I&0&D&-D\\S&D&-D&D&0&-I\\S&D&D&-D&-I&0
\end{pmatrix},\qquad W=A(I_6\otimes F).
\]
As certified previously, $WW^{\mathsf T}=10I_{12}$, but this input
admits no zero-preserving Hadamard filling. The checker re-verifies
that obstruction by exhausting all eight root states.

Apply the theorem's $Q$ to rows $0,1,2,3$, with all indices zero-based.
Every outside-column product is $-1$. The first four rows of the
switched weighing matrix $\widetilde W$ become
\[
\begin{pmatrix}
1&0&1&0&-1&1&-1&1&1&1&1&1\\
0&-1&0&1&1&-1&1&-1&1&1&1&1\\
-1&0&1&0&1&1&1&1&-1&1&-1&1\\
0&1&0&1&1&1&1&1&1&-1&1&-1
\end{pmatrix};
\]
the other eight rows are unchanged. In particular the two original
zero-row pairs are reorganized into pairs $(0,2)$ and $(1,3)$.

Here is a complete filling certificate. In each table row insert
the indicated signs in the two listed zero columns, leaving all
nonzeros of $\widetilde W$ unchanged:
\[
\begin{array}{c|c|c}
\text{row}&\text{zero columns}&\text{inserted signs}\\\hline
0&(1,3)&(-1,-1)\\1&(0,2)&(-1,1)\\
2&(1,3)&(1,-1)\\3&(0,2)&(1,1)\\
4&(4,5)&(1,-1)\\5&(4,5)&(-1,-1)\\
6&(6,7)&(1,-1)\\7&(6,7)&(-1,-1)\\
8&(8,9)&(-1,1)\\9&(8,9)&(1,1)\\
10&(10,11)&(-1,1)\\11&(10,11)&(1,1)
\end{array}
\]
The resulting sign matrix has exact Gram $12I_{12}$. Negating
all inserted signs gives the second filling; the eight-state test
certifies that these are the only two fillings of $\widetilde W$.

The switch changes twenty-six matrix entries. Twenty-two originally
nonzero entries change or become zero, and eight positions change
zero/nonzero status (four holes are erased and four are introduced).
Each completed Hadamard matrix disagrees with only twenty of the
original $W$'s 120 prescribed nonzero signs; all disagreements lie
in the four switched rows. These are exact measured costs, not
proved global minima. For this fixed orientation of $Q$, precisely
one of the fifteen original row-pair choices is admissible.

The checker also records a bounded failed alternative: among 3,200
specified equal-magnitude Householder selectors, none produced a
nontrivial zero-preserving weighing switch on this input. This does
not assert isolation under arbitrary operations. The successful
four-row move preserves the weighing equation while changing the
zero pairing, which the fixed-support rigidity proposition forbids
for a feasibility-changing move of the same row-block size.

This supplies an input-producing operation and an actual repair
witness. It neither guarantees that a suitable quadruple exists
for every weighing matrix nor shows that iterating these switches
will always reach a completable input.

\section{The remaining construction problem}
Continuous phase escape is available at every non-Hadamard real sign
matrix, but a phase trajectory can approach a complex Hadamard matrix
outside the four-phase alphabet. No theorem makes phase dispersion
decrease along the new direction. The completion algorithm needs all
but a bounded number of columns in advance, and weighing exchanges
must preserve both their algebraic equations and their zero support.
None of these hypotheses is supplied uniformly at unknown orders.
The next target is an operation coupling orthogonality improvement to
phase quantization, or an input-producing theorem for the bounded
completion algorithm. No Lean edits, git commits or submissions were made.
\end{document}
